\documentclass[a4paper, serif, 10pt]{moderncv}

%% moderncv
% style and color
\moderncvstyle{banking}
\moderncvcolor{black}

% renew moderncv command to adapt for CJK colon
\renewcommand*{\cvitem}[3][.25em]{%
  \ifstrempty{#2}{}{\hintstyle{#2}：}{#3}%
  \par\addvspace{#1}}

\renewcommand*{\cvitemwithcomment}[4][.25em]{%
  \savebox{\cvitemwithcommentmainbox}{\ifstrempty{#2}{}{\hintstyle{#2}：}#3}%
  \setlength{\cvitemwithcommentmainlength}{\widthof{\usebox{\cvitemwithcommentmainbox}}}%
  \setlength{\cvitemwithcommentcommentlength}{\maincolumnwidth-\separatorcolumnwidth-\cvitemwithcommentmainlength}%
  \begin{minipage}[t]{\cvitemwithcommentmainlength}\usebox{\cvitemwithcommentmainbox}\end{minipage}%
  \hfill% fill of \separatorcolumnwidth
  \begin{minipage}[t]{\cvitemwithcommentcommentlength}\raggedleft\small\itshape#4\end{minipage}%
  \par\addvspace{#1}}

%% page layout/margins
\usepackage[top=2.5cm, bottom=2.5cm, left=1.5cm, right=1.5cm]{geometry}

\nopagenumbers{}

%% Babel config for English language
\usepackage[english]{babel}

%% fontspec
\usepackage{fontspec}

\IfFontExistsTF{Linux Libertine}{
  \setmainfont[Ligatures={TeX, Common}, Numbers=Lining, ItalicFont=Linux Libertine]{Linux Libertine}
}{}
\IfFontExistsTF{Linux Libertine O}{
  \setmainfont[Ligatures={TeX, Common}, Numbers=Lining, ItalicFont=Linux Libertine O]{Linux Libertine O}
}{}

%% CTeX
% CJK support, used to show CJK characters in the resume
%
% - fontset=none: disable builtin fontset but instead set the CJK font manually
% - heading=false: disable ctex heading
% - punct=kaiming: use kaiming punctuations styles for CJK
% - scheme=plain: use plain scheme, do not override \normalsize font size
% - space=auto: space settings for CJK characters
%
% ref:
% - http://ctan.mirrorcatalogs.com/language/chinese/ctex/ctex.pdf
\usepackage[UTF8, heading=false, punct=kaiming, scheme=plain, space=auto]{ctex}

\IfFontExistsTF{Noto Serif CJK SC}{
  \setCJKmainfont{Noto Serif CJK SC}
}{}
\IfFontExistsTF{Noto Sans CJK SC}{
  \setCJKsansfont{Noto Sans CJK SC}
}{}

%% line spacing
\usepackage{setspace}
\setstretch{1.125}

%% URL styling - use normal text instead of monospace
\urlstyle{same}

% Auto-underline all links
% Defer until \begin{document} so hyperref is already loaded
\AtBeginDocument{%
  \let\oldhref\href
  \renewcommand{\href}[2]{\oldhref{#1}{\underline{#2}}}%
}

%% Patch \httplink and \httpslink to handle full URLs
% The original moderncv macros always prepend http:// or https:// to URLs.
% This causes issues when the URL already has a protocol (e.g., https://example.com)
% resulting in malformed URLs like https://https://example.com.
% This patch checks if the URL already contains "://" and uses it directly if so.
% Note: We use \str_set:Nx (x-expansion) to expand macros like \@homepage first.
\ExplSyntaxOn
\str_new:N \g_yamlresume_url_str

\RenewDocumentCommand{\httpslink}{O{}m}{
  \str_set:Nx \g_yamlresume_url_str {#2}
  \str_if_in:NnTF \g_yamlresume_url_str {://}
    { \url{#2} }
    { \url{https://#2} }
}

\RenewDocumentCommand{\httplink}{O{}m}{
  \str_set:Nx \g_yamlresume_url_str {#2}
  \str_if_in:NnTF \g_yamlresume_url_str {://}
    { \url{#2} }
    { \url{http://#2} }
}
\ExplSyntaxOff

\name{佐藤美咲}{}
\title{データアナリスト | データ可視化・ビジネスインサイト}
\phone[mobile]{080-1234-5678}
\email{misaki.sato@example.com}
\homepage{https://misaki-sato.dev}

\address{日本、渋谷区、東京都、東京都渋谷区神南1-2-3、150-0001}{}{}

\extrainfo{{\small \faGithub}\ \href{https://github.com/misaki-sato}{@misaki-sato} {} {} {} • {} {} {}
{\small \faLinkedin}\ \href{https://www.linkedin.com/in/misaki-sato}{@misaki-sato}}

\begin{document}

\maketitle

\section{基本情報}

\cvline{}{データに基づく意思決定を支援するデータアナリスト。小売・EC業界でSQL・Python・Tableauを駆使し、売上分析、顧客セグメンテーション、予測モデル構築などに従事。ビジネス課題を数値で可視化し、チームのKPI改善に貢献。統計・機械学習の知識を活かした実践的な分析に強みを持つ。}
\section{学歴}

\cventry{2017年4月～2021年3月}
        {学士、経済学部 経済学科、成績：3.7}
        {東京大学}
        {\url{https://www.u-tokyo.ac.jp/}}
        {}
        {\begin{itemize}
\item 統計学・計量経済学を中心に、データ分析の基礎を学ぶ
\item 卒業論文では、POSデータを用いた購買行動分析を実施
\item 学外のデータ分析コンテストに参加し、可視化とプレゼンテーションを経験
\end{itemize}
\textbf{コース}：統計学、計量経済学、データ解析基礎、プログラミング入門 (Python)、マーケティングリサーチ}
\section{職歴}

\cventry{2021年4月～現在}
        {データアナリスト}
        {株式会社インサイト}
        {\url{https://www.insight-example.co.jp}}
        {}
        {\begin{itemize}
\item 小売・EC事業の売上・在庫・顧客データを統合し、ダッシュボードを構築
\item SQL・Pythonによるデータ抽出・加工・分析を行い、週次・月次レポートを自動作成
\item 顧客セグメンテーション分析を実施し、マーケティング施策の効果検証を支援
\item Tableauを用いて経営陣向けのKPIダッシュボードを開発し、意思決定を支援
\end{itemize}
\textbf{キーワード}：SQL、Python、Tableau、KPI分析、ダッシュボード}

\cventry{2020年7月～2021年3月}
        {リサーチアシスタント}
        {株式会社ブライトリサーチ}
        {\url{https://www.bright-research-example.co.jp}}
        {}
        {\begin{itemize}
\item 市場調査データのクリーニング・集計を担当
\item Excel・Rを用いた統計解析と結果の可視化を支援
\item 調査レポートの作成に協力し、クライアントへの提案を補助
\end{itemize}
\textbf{キーワード}：Excel、R、統計解析、市場調査}
\section{言語}

\cvline{日本語}{ネイティブ・母国語レベル \hfill \textbf{キーワード}：母国語}
\cvline{英語}{完全な実務レベル \hfill \textbf{キーワード}：TOEIC 920、ビジネス文書作成}
\section{スキル}

\cvline{データ分析}{エキスパート \hfill \textbf{キーワード}：Python、R、SQL、統計解析、予測モデル}
\cvline{データ可視化}{上級 \hfill \textbf{キーワード}：Tableau、Power BI、Looker Studio、matplotlib、ggplot2}
\cvline{ビジネスインテリジェンス}{上級 \hfill \textbf{キーワード}：KPI管理、マーケティング分析、顧客セグメンテーション、A/Bテスト}
\cvline{機械学習}{中級 \hfill \textbf{キーワード}：scikit-learn、回帰分析、クラスタリング、意思決定木}
\section{受賞・表彰}

\cventry{2022年12月}
        {株式会社インサイト}
        {社内データ分析コンテスト 最優秀賞}
        {}
        {}
        {購買データを用いた顧客離反予測モデルを提案し、社内のビジネスコンテストで最優秀賞を受賞。}
\section{資格・免状}

\cventry{2022年6月}
        {日本統計学会}
        {統計検定2級}
        {\url{https://www.toukei-kentei.jp/}}
        {}
        {}

\cventry{2021年3月}
        {Google}
        {Google Data Analytics Professional Certificate}
        {\url{https://www.coursera.org/professional-certificates/google-data-analytics}}
        {}
        {}
\section{出版・著作}

\cventry{2023年11月}
        {データドリブンマーケティングの実践: 小売業における顧客セグメント分析}
        {社内テクニカルレポート}
        {\url{https://misaki-sato.dev/articles/data-driven-marketing}}
        {}
        {小売業の購買データを活用した顧客セグメンテーションの実践手法をまとめた。RFM分析とクラスタリングを組み合わせたアプローチを紹介。}
\section{プロジェクト}

\cventry{2023年1月～2023年6月}
        {ECサイトの売上予測モデルと可視化ダッシュボードを開発。}
        {EC売上予測ダッシュボード}
        {\url{https://misaki-sato.dev/projects/ec-sales-forecast}}
        {}
        {\begin{itemize}
\item 過去3年分の売上データを基に、季節性を考慮した予測モデルを構築
\item StreamlitとPlotlyを用いて、担当者が簡単に予測を確認できるダッシュボードを提供
\item 在庫計画やプロモーション効果の事前検証に貢献
\end{itemize}
\textbf{キーワード}：Python、Streamlit、Plotly、時系列分析}

\cventry{2022年4月～2022年9月}
        {購買履歴を基にした顧客セグメンテーションの分析プロジェクト。}
        {顧客セグメンテーション分析}
        {\url{https://misaki-sato.dev/projects/segmentation}}
        {}
        {\begin{itemize}
\item RFM分析とk-meansクラスタリングを組み合わせて顧客を分類
\item セグメントごとの特徴を分析し、マーケティング施策の提案につなげた
\item リピート率向上に向けた施策の優先順位付けを支援
\end{itemize}
\textbf{キーワード}：RFM分析、k-means、Tableau、マーケティング}
\section{興味・関心}

\cvline{ランニング}{マラソン、トレーニング}
\cvline{読書}{データサイエンス、ビジネス書}
\cvline{データ可視化コミュニティ}{勉強会、ハンズオン}
\section{ボランティア活動}

\cventry{2022年4月～現在}
        {データ分析ボランティア}
        {NPO法人 データ支援ネット}
        {\url{https://www.data-support-npo.example.org}}
        {}
        {\begin{itemize}
\item 非営利団体向けに、アンケート集計や業務改善のためのデータ分析を支援
\item 事務局メンバー向けにExcelや可視化ツールの使い方研修を実施
\end{itemize}}

\end{document}